\begin{statementbox}
\begin{center}\begin{tikzpicture}[font=\small,>=stealth,line join=round]
\draw[gray!40,rounded corners=3pt] (0,0) rectangle (14.2,8.3);\node at (7.1,7.700000000000001) {两条曲线都过(0,1)：向右走，一条升、一条降};
\begin{scope}[shift={(4.1,2.068181818181818)},x=1.0714285714285714cm,y=0.6969696969696969cm]\draw[->,gray] (-2.8,0)--(2.8,0) node[right]{$x$};\draw[->,gray] (0,-0.6)--(0,6) node[above]{$y$};\draw[blue!75!black,thick,solid,domain=-2.5:2.15,samples=180] plot(\x,{pow(2,\x)});\draw[orange!80!black,thick,dashed,domain=-2.5:2.15,samples=180] plot(\x,{pow(2,-\x)});\fill[orange!85!black] (0,1) circle(2.2pt);\fill[orange!85!black] (1,2) circle(2.2pt);\fill[orange!85!black] (-1,2) circle(2.2pt);\node[below=5pt] at (-2,0){$-2$};\node[below left=5pt] at (0,0){$0$};\node[below=5pt] at (2,0){$2$};\draw[->,blue!75!black,thick] (.6,1.5157)--(1.2,2.2974);\draw[->,orange!80!black,thick] (.7,.6156)--(1.5,.3536);\end{scope}\node[anchor=west,align=left,text width=5.0cm] at (8.4,4.0) {蓝实线：$y=2^x$，向右升\\[9pt]橙虚线：$y=(\frac12)^x$，向右降\\[9pt]共同点：$(0,1)$\\[9pt]整条曲线都在横轴上方\\[9pt]比较时先看底数与$1$};
\node[text width=12.6cm,align=center] at (7.1,.65) {零次方都等于1；每向右一步，蓝线高度乘2，橙线高度乘二分之一。};
\end{tikzpicture}\end{center}

\textbf{观察：}看共同点、横轴上方和两条曲线的方向。\quad\textbf{原因：}零次方为1，正底数的幂为正。

\textbf{考试使用：}认图先找(0,1)，再按底数判断增减。



\textbf{\textcolor{CStatement}{条件}}
\begin{description}[style=nextline, leftmargin=2.8em, labelsep=0.5em, itemsep=0.35em]
    \item[\textbf{\textcolor{CStatement}{底数与定义域：}}]
    设 $f(x)=a^x$，其中 $a>0$、$a\neq1$，$x\in\mathbb R$。这里是\textbf{底数 $a$ 固定，指数 $x$ 变化}；不要与幂函数 $y=x^a$ 混淆。
\end{description}

\textbf{\textcolor{CStatement}{结论}}
\begin{description}[style=nextline, leftmargin=2.8em, labelsep=0.5em, itemsep=0.45em]
    \item[\textbf{\textcolor{CStatement}{结论一（正值、定点、值域与渐近线）：}}]
    对任何允许的底数，$a^0=1$ 且 $a^x>0$；图像总经过 $(0,1)$，始终在 $x$ 轴上方，能无限接近 $x$ 轴但不会碰到它。
    \[
        \boxed{\ f(0)=1,\qquad a^x>0,\qquad \text{值域 }(0,+\infty),\qquad y=0\text{ 为水平渐近线}\ }
    \]
    当 $a>1$ 时，向左走得越远，函数值越接近 $0$；当 $0<a<1$ 时，向右走得越远，函数值越接近 $0$。

\end{description}
\end{statementbox}
\clearpage
\begin{statementbox}
\begin{description}[style=nextline, leftmargin=2.8em, labelsep=0.5em, itemsep=0.45em]
    \item[\textbf{\textcolor{CStatement}{结论二（增减由底数决定）：}}]
    \[
        \begin{array}{c|c|c|c}
        \text{底数} & x\text{ 从小到大} & x<0 & x>0\\ \hline
        a>1 & a^x\text{ 递增} & 0<a^x<1 & a^x>1\\
        0<a<1 & a^x\text{ 递减} & a^x>1 & 0<a^x<1
        \end{array}
    \]
    也就是说，若 $x_1<x_2$，则
    \[
        \begin{cases}
        a^{x_1}<a^{x_2},& a>1,\\
        a^{x_1}>a^{x_2},& 0<a<1.
        \end{cases}
    \]
    \textbf{理解抓手：}指数 $x$ 每增加 $1$，函数值就\textbf{乘一次 $a$}：$a^{x+1}=a\cdot a^x$。乘大于 $1$ 的数会变大；乘 $0$ 到 $1$ 之间的数会变小。
\end{description}

\textbf{\textcolor{CStatement}{推广结论（两条图像关于 $y$ 轴对称）}}
\begin{description}[style=nextline, leftmargin=2.8em, labelsep=0.5em, itemsep=0.45em]
    \item[\textbf{\textcolor{CStatement}{底数互为倒数：}}]
    令 $g(x)=\left(\frac1a\right)^x$，有
    \[
        \boxed{\ g(x)=a^{-x}=f(-x)\ }.
    \]
    因而 $y=a^x$ 与 $y=(1/a)^x$ 关于 $y$ 轴对称：前者图像上的点 $(t,a^t)$ 对应后者图像上的点 $(-t,a^t)$。\textbf{横坐标变号，纵坐标不变。}
\end{description}

\end{statementbox}
